% ===== §2 =====
\section{The Formal Anatomy of an Optimization Problem}
\label{sec:anatomy}

\noindent
Before one can show that intimate relation furnishes no well-defined optimization problem, one must be exact about what a well-defined optimization problem requires, and the exactness is not pedantry but the whole condition of the argument's force. The four impossibilities of the paper's central part are, each of them, the failure of one of the conditions set out here; if the conditions are stated loosely, the failures can be dodged, and if they are stated in a caricatured form the whole refutation is a straw man. The aim of this section is therefore to state, without polemic and in a form a theorist of decision would himself accept, precisely what has to be in place for the operation of maximization to denote anything at all. The section grants optimization its strongest and most respectable version, so that when the later parts withdraw the conditions one by one, they withdraw them from a construction the optimizer recognizes as his own and not from a parody built to be knocked down.

An optimization problem, in its canonical form, is the search for an element of a set that makes a function as large as possible. Written out, it is the expression
\[
x^{\ast} \;=\; \argmax_{x \in \mathcal{X}} \; J(x),
\]
and the compactness of the notation hides four presuppositions, each of which must hold for the expression to pick out anything. There is a \emph{feasible set} $\mathcal{X}$, the space of available options, which must be given in advance and be well defined, so that it is determinate what is and is not a candidate. And there is an \emph{objective function} $J$, which assigns to each candidate a value; and it is in the conditions on $J$ that the presuppositions grow heavy. For the expression to denote, $J$ must be \emph{representationally complete}, defined over terms into which every value at stake has been written, since a value absent from $J$ is a value the maximization cannot see. It must be \emph{exogenous and fixed}, given prior to the search and unchanged by it, for a function that drifts while one climbs offers no summit one can reach and keep. It must be \emph{scalar}, or valued in an ordering rich enough that maxima exist, since $\argmax$ over values that do not all compare has nothing determinate to return. And behind all three there must be an \emph{optimizer}: a subject whose objective $J$ is, whose preferences it encodes, presupposed to stand outside the problem, prior to it and stable through it, as the one for whom $x^{\ast}$ is best.

\subsection*{The four presuppositions, named in the order of their denial}

It is worth naming these four in the order the paper will deny them, since the order is a descent from the surface of the problem to its ground, from whether the problem can be \emph{written} to who is doing the \emph{writing}, and each later condition presupposes the earlier ones already met.

The first presupposition is \emph{representational completeness}: that every value at stake can be set down as a term in $J$. This is the condition that the problem can be stated at all, logically prior to every other, for a value that cannot enter $J$ is not priced badly within the optimization but excluded from its statement, beyond the reach of any adjustment of the function. The first impossibility denies this.

The second is \emph{exogeneity and fixity}: that the objective is given before the process and unaltered by it, a determinate target rather than a moving one. This is the condition an endogenous, self-revising objective violates, and the second impossibility denies it, not as a contingent fact about people whose tastes happen to wander but on the ontological ground that value in relation is generated within the relation and admits no external position from which it could be fixed.

The third is \emph{scalarity}: that the values $J$ assigns lie on a common scale, so that any two candidates compare and $\argmax$ has something to select. This is the condition incommensurable goods violate, and the third impossibility denies it.

The fourth is the \emph{stable optimizer}: a subject with settled preferences standing outside the problem, whose objective $J$ is and for whom its maximum is best. This is the condition a subject constituted within the very relation being optimized violates, and the fourth impossibility denies it, dissolving the fixed referent of the ``who'' the whole construction silently presupposes.

\subsection*{An instrument granted its proper domain}

It clears the ground to see the four conditions holding somewhere, for the paper does not claim that optimization is everywhere ill-posed, only that it is so in a particular and important region, and the contrast is sharper if the well-posed case is granted its due. Consider a shipping company routing a fleet to minimize fuel cost. The feasible set is the routes the vehicles can take, well defined by the road network; the objective is fuel burned, representationally complete because fuel is the whole of what is being weighed and nothing relevant is left outside it; exogenous and fixed, since the price of fuel and the shape of the network do not shift because the planner is planning; scalar, since fuel is measured in one unit along which more and less are always defined; and owned by a stable optimizer, the firm, whose objective this determinately is and who is not remade by the routing. Here optimization is not merely legitimate but the only serious way to proceed, and the paper affirms this without reservation. The four conditions are jointly met, the operation is well defined, and to decline it would be a failure of reason and not a fidelity to anything. The example matters precisely because it is so clean: it shows that the paper's quarrel is not with optimization but with its transposition into a domain where, as the next part argues, not one of these four conditions survives.

\subsection*{The demand that the anatomy be granted in full}

The strategy of the argument requires that this anatomy be granted in its full strength, and the reason is worth stating, because it marks off the paper's claim from a familiar and much weaker one with which it would otherwise be confused. The weaker claim holds that optimization problems in human affairs are merely \emph{hard}: that the feasible set is too vast to search, the objective too costly to evaluate, the information too poor, so that we fall back of necessity on rules of thumb and on stopping when a candidate is good enough rather than best \citep{simon1955behavioral}. This is the objection from bounded rationality, and it is true and important, and it is not this paper's objection; indeed it concedes exactly what this paper denies. To satisfice because the true optimum is out of computational reach is to grant that a true optimum exists, well defined but unreachable, and merely to humble our access to it; the metaphysics of optimization is left standing, and only our finitude is confessed. Even the theory of choice under uncertainty, which replaces the known objective with an expectation over states, preserves the four conditions intact and simply integrates $J$ against a probability, so that the agent maximizes expected rather than certain utility and the whole apparatus, with its complete, fixed, scalar objective and its stable subject, is carried over undisturbed \citep{savage1954foundations}. The paper's claim reaches beneath all of this. It is not that the problem is well-posed and hard, nor that it is well-posed and risky, but that it is not well-posed at all, that one or more of the four conditions fails in kind and not in degree. To establish that, one must take optimization at its strongest, grant it a searcher of unlimited power and perfect information facing no uncertainty whatever, and show that even for such a searcher the problem does not exist, because the representationally complete, exogenous, fixed, scalar objective belonging to a stable subject that the operation requires is, in intimate relation, not there to be operated on. The anatomy is set out in full so that the denials, when they come, cannot be answered by the reply that a faster computer or a more patient economist would have found the maximum after all. There is, the next part argues, no maximum there to be found, and no faster computer reaches what was never posed.
